\documentclass[11pt,a4paper]{article}
\usepackage[utf8]{inputenc}
\usepackage[T1]{fontenc}
\usepackage{amsfonts}
\usepackage[french]{babel}
\frenchbsetup{StandardLists=true}
\usepackage{enumitem}
\usepackage{amssymb}
\usepackage{amsmath}

\usepackage[left=2cm,right=2cm,top=2cm,bottom=2cm,headheight=14pt]{geometry}
\usepackage{multirow}
\usepackage[dvipsnames]{xcolor}
\usepackage{parskip}
\usepackage{caption}

\usepackage[T1]{fontenc}
\usepackage{fouriernc}

\usepackage{graphicx}
\usepackage{setspace}
\singlespacing

\usepackage{tcolorbox}
\usepackage{framed}
\usepackage{multicol}

\usepackage{fancybox}
\usepackage{fancyhdr}
\pagestyle{fancy}
\renewcommand\headrulewidth{1pt}
\fancyhead[L]{Lycée Denis Diderot}
\fancyhead[R]{Première Sciences}
\setcounter{page}{1}\fancyfoot[C]{TP Noté -- Loi de Boyle-Mariotte -- Page \thepage}
\fancyfoot[R]{Année 2024 - 2025}
\fancyfoot[L]{Alexis RUIZ}

\usepackage{pstricks,pst-plot}
\usepackage{tikz}
\usepackage{tkz-tab}
\usetikzlibrary{arrows,calc}
\usepackage{pifont}

\usepackage[tikz]{bclogo}

% Commande pour les lignes de réponse
\newcommand{\lignesreponse}[1][3]{%
\vspace{-2mm}
    \foreach \n in {1,...,#1}{%
        \noindent\makebox[\linewidth]{\dotfill}\\[0.5cm]
    }
}

% Boîte barème (points)
\newcommand{\bareme}[1]{%
    \hfill\colorbox{gray!20}{\small\textbf{/ #1 pt\ifnum#1>1 s\fi}}%
}

% Boîtes colorées
\newtcolorbox{contextebox}{
    colback=blue!5,
    colframe=blue!50!black,
    title=Contexte,
    fonttitle=\bfseries,
    arc=3mm,
    boxrule=1pt
}

\newtcolorbox{documentbox}[1]{
    colback=green!5,
    colframe=green!50!black,
    title=Document #1,
    fonttitle=\bfseries,
    arc=3mm,
    boxrule=1pt
}

\newtcolorbox{rappelbox}{
    colback=yellow!10,
    colframe=orange!75!black,
    title=Rappel,
    fonttitle=\bfseries,
    arc=3mm,
    boxrule=1pt
}

\newtcolorbox{baremeBox}{
    colback=gray!10,
    colframe=gray!60!black,
    title=Barème,
    fonttitle=\bfseries,
    arc=3mm,
    boxrule=1pt
}

\newtcolorbox{attentionbox}{
    colback=orange!5,
    colframe=orange!75!black,
    title=Attention,
    fonttitle=\bfseries,
    arc=3mm,
    boxrule=1pt
}

\begin{document}

%-------------------------------------------------------------------------------
% EN-TÊTE ÉLÈVE
%-------------------------------------------------------------------------------
\begin{center}
 \LARGE{\textbf{TP Noté -- Loi de Boyle-Mariotte}}

 \large{Pression et comportement des gaz}
\end{center}

\vspace{3mm}
\hrule
\vspace{4mm}

\begin{minipage}{0.55\linewidth}
    \textbf{Nom :} \dotfill \\[6pt]
    \textbf{Prénom :} \dotfill \\[6pt]
    \textbf{Classe :} \dotfill \hspace{1cm} \textbf{Date :} \dotfill
\end{minipage}
\hfill
\begin{minipage}{0.38\linewidth}
    \begin{baremeBox}
    \begin{center}
    \begin{tabular}{lr}
        Partie 1 & \quad /4 pts \\
        Partie 2 & \quad /6 pts \\
        Partie 3 & \quad /6 pts \\
        Partie 4 & \quad /4 pts \\
        \hline
        \textbf{Total} & \textbf{\quad /20 pts} \\
    \end{tabular}
    \end{center}
    \end{baremeBox}
\end{minipage}

\vspace{5mm}
\hrule
\vspace{5mm}

%-------------------------------------------------------------------------------
% CONTEXTE
%-------------------------------------------------------------------------------

\begin{contextebox}
Un technicien d'un service médical souhaite étudier le comportement de l'air comprimé dans une seringue. Pour cela, il a réalisé une expérience en connectant une seringue graduée à un capteur de pression MPX5700AP relié à un microcontrôleur Arduino. Il a enregistré des couples de valeurs (pression, volume) en faisant varier la position du piston, à \textbf{température constante}.

\vspace{2mm}
\textbf{Problématique :} Comment varie la pression d'un gaz lorsque son volume change à température constante ?
\end{contextebox}

\vspace{3mm}

%-------------------------------------------------------------------------------
% DOCUMENTS
%-------------------------------------------------------------------------------

\begin{documentbox}{1 -- Données expérimentales recueillies}
Le technicien a relevé les valeurs suivantes lors de son expérience :

\vspace{3mm}
\begin{center}
\begin{tabular}{|c|c|c|c|c|c|c|c|}
\hline
\textbf{Volume V (mL)} & 20 & 25 & 30 & 35 & 40 & 50 & 60 \\
\hline
\textbf{Pression P (kPa)} & 201,5 & 162,1 & 133,8 & 115,9 & 101,2 & 80,5 & 67,6 \\
\hline
\textbf{$\dfrac{1}{V}$ (mL$^{-1}$)} & \hspace{1cm} & \hspace{1cm} & \hspace{1cm} & \hspace{1cm} & \hspace{1cm} & \hspace{1cm} & \hspace{1cm} \\[6pt]
\hline
\textbf{P $\times$ V (kPa·mL)} & \hspace{1cm} & \hspace{1cm} & \hspace{1cm} & \hspace{1cm} & \hspace{1cm} & \hspace{1cm} & \hspace{1cm} \\[6pt]
\hline
\end{tabular}
\end{center}
\end{documentbox}

\vspace{3mm}

\begin{documentbox}{2 -- Graphique obtenu par le programme Python}
Le programme Python a tracé le graphique de la pression $P$ en fonction de $\dfrac{1}{V}$. La droite de régression linéaire obtenue a pour équation :

$$P = 4045 \times \dfrac{1}{V} - 0{,}9 \quad \text{(avec P en kPa et V en mL)}$$

Le coefficient de corrélation vaut : $R^2 = 0{,}9998$

\vspace{3mm}

% Graphique tikz représentant P = f(1/V)
\begin{center}
\begin{tikzpicture}
\begin{scope}[xscale=90, yscale=0.040]
    % Axes
    \draw[->] (0,0) -- (0.057,0) node[right] {$\dfrac{1}{V}$ (mL$^{-1}$)};
    \draw[->] (0,0) -- (0,215) node[above] {$P$ (kPa)};
    % Graduation axe X
    \foreach \x/\xl in {0.010/0.010, 0.020/0.020, 0.030/0.030, 0.040/0.040, 0.050/0.050}{
        \draw (\x, 2pt/0.040) -- (\x, -2pt/0.040) node[below, font=\small] {\xl};
    }
    % Graduation axe Y
    \foreach \y in {50, 100, 150, 200}{
        \draw (0.001, \y) -- (-0.001, \y) node[left, font=\small] {\y};
    }
    % Droite de régression : P = 4045/V - 0.9
    \draw[blue, thick] (0.0165, 4045*0.0165-0.9) -- (0.052, 4045*0.052-0.9);
    % Points expérimentaux
    \foreach \v/\p in {20/201.5, 25/162.1, 30/133.8, 35/115.9, 40/101.2, 50/80.5, 60/67.6}{
        \pgfmathsetmacro{\xinv}{1/\v}
        \fill[red] (\xinv, \p) circle (1pt);
    }
    % Légende
    \draw[blue, thick] (0.035, 170) -- (0.042, 170) node[right, font=\small] {Régression linéaire};
    \fill[red] (0.0375, 158) circle (1pt) node[right, font=\small] {Points mesurés};
\end{scope}
\end{tikzpicture}
\end{center}
\end{documentbox}

\vspace{3mm}

\begin{rappelbox}
\textbf{Loi de Boyle-Mariotte :} À température constante et pour une quantité de matière donnée, le produit de la pression $P$ par le volume $V$ d'un gaz est constant :
$$\boxed{P \times V = \text{constante}} \quad \Longleftrightarrow \quad P_1 \times V_1 = P_2 \times V_2$$

\textbf{Unités de pression courantes :}
\begin{itemize}[leftmargin=*]
    \item $1\,\text{bar} = 10^5\,\text{Pa} = 100\,\text{kPa}$
    \item $1\,\text{hPa} = 100\,\text{Pa}$
    \item Pression atmosphérique : $P_{atm} \approx 1013\,\text{hPa} \approx 101{,}3\,\text{kPa}$
\end{itemize}
\end{rappelbox}

\newpage

%-------------------------------------------------------------------------------
% PARTIE 1 -- CONNAISSANCES
%-------------------------------------------------------------------------------

\section*{Partie 1 -- Connaissances sur la pression \bareme{4}}

\begin{enumerate}

    \item \textbf{Donner} l'unité internationale (SI) de la pression et son symbole. \bareme{1}

    \lignesreponse[1]

    \item \textbf{Convertir} la pression atmosphérique $P_{atm} = 1013\,\text{hPa}$ en kPa et en Pa. \bareme{1}

    \vspace{2mm}
    Calcul : \dotfill

    \vspace{3mm}
    $P_{atm} = $ \dotfill kPa $= $ \dotfill Pa

    \vspace{3mm}

    \item \textbf{Énoncer} la loi de Boyle-Mariotte avec vos propres mots (sans formule). \bareme{1}

    \lignesreponse[2]

    \item \textbf{Selon} la loi de Boyle-Mariotte, si on réduit le volume d'un gaz de moitié (à température constante), que devient sa pression ? Justifier. \bareme{1}

    \lignesreponse[2]

\end{enumerate}

%-------------------------------------------------------------------------------
% PARTIE 2 -- EXPLOITATION DU TABLEAU
%-------------------------------------------------------------------------------

\section*{Partie 2 -- Exploitation du tableau de données \bareme{6}}

\begin{enumerate}

    \item \textbf{Calculer} et \textbf{compléter} les valeurs de $\dfrac{1}{V}$ dans le tableau du Document 1. Arrondir à $10^{-3}$ près. \bareme{2}

    \vspace{2mm}
    \textit{(Inscrire les valeurs directement dans le tableau.)}

    \vspace{3mm}

    \item \textbf{Calculer} et \textbf{compléter} les valeurs de $P \times V$ dans le tableau du Document 1. Arrondir à l'unité. \bareme{2}

    \vspace{2mm}
    Montrer un exemple de calcul ci-dessous :

    \vspace{5mm}
    Calcul pour $V = 20\,\text{mL}$ : \dotfill

    \vspace{5mm}

    \item Les valeurs de $P \times V$ sont-elles constantes ? \textbf{Calculer} la valeur moyenne de $P \times V$ et \textbf{conclure}. \bareme{2}

    \vspace{5mm}
    Calcul de la moyenne : \dotfill

    \vspace{5mm}
    Moyenne de $P \times V = $ \dotfill kPa·mL

    \vspace{3mm}

    \lignesreponse[2]

\end{enumerate}

%-------------------------------------------------------------------------------
% PARTIE 3 -- ANALYSE GRAPHIQUE
%-------------------------------------------------------------------------------

\section*{Partie 3 -- Analyse du graphique \bareme{6}}

\begin{enumerate}

    \item \textbf{Décrire} l'allure de la courbe $P = f\!\left(\dfrac{1}{V}\right)$ visible sur le graphique du Document 2. \bareme{1}

    \lignesreponse[1]

    \item \textbf{Relever} les paramètres de la droite de régression donnée dans le Document 2 : \bareme{1}

    \vspace{2mm}
    Coefficient directeur : $a = $ \dotfill

    \vspace{3mm}
    Ordonnée à l'origine : $b = $ \dotfill

    \vspace{3mm}
    Coefficient de corrélation : $R^2 = $ \dotfill

    \vspace{3mm}

    \item Le coefficient de corrélation $R^2 = 0{,}9998$. Qu'indique cette valeur sur la qualité de la régression linéaire ? \bareme{1}

    \lignesreponse[1]

    \item La valeur de $b$ est très proche de zéro. \textbf{Quelle relation} cela implique-t-il entre $P$ et $\dfrac{1}{V}$ ? \bareme{1}

    \lignesreponse[1]

    \item \textbf{En déduire} la relation entre $P$ et $V$ (en utilisant $a$ comme constante). \bareme{1}

    \vspace{5mm}

    On peut écrire : $P = $ \dotfill, soit $P \times V = $ \dotfill

    \vspace{5mm}

    \item \textbf{Comparer} la valeur du coefficient directeur $a$ avec la moyenne de $P \times V$ calculée en Partie 2. Que remarque-t-on ? \bareme{1}

    \lignesreponse[2]

\end{enumerate}

%-------------------------------------------------------------------------------
% PARTIE 4 -- CONCLUSION
%-------------------------------------------------------------------------------

\section*{Partie 4 -- Conclusion \bareme{4}}

\begin{enumerate}

    \item Les résultats obtenus sont-ils en accord avec la loi de Boyle-Mariotte ? \textbf{Justifier} en vous appuyant sur les valeurs numériques des parties précédentes. \bareme{2}

    \lignesreponse[3]

    \item \textbf{Citer} une source d'incertitude expérimentale pouvant expliquer que les produits $P \times V$ ne soient pas parfaitement constants. \bareme{1}

    \lignesreponse[2]

    \item \textbf{Proposer} une application concrète de la loi de Boyle-Mariotte dans la vie quotidienne ou dans un domaine technique. \bareme{1}

    \lignesreponse[2]

\end{enumerate}

\vspace{5mm}
\hrule
\vspace{3mm}
\begin{center}
\textit{--- Fin du TP noté ---}
\end{center}

\end{document}
